\begin{theorem}[$S = 1/\sqrt{\text{PP}}$]
  \label{thm:equivalence}
  For any polygon $D$ with area $A$ and perimeter $P$:
  \begin{equation}
    S(D) = \frac{1}{\sqrt{\text{PP}(D)}}.
    \label{eq:equivalence}
  \end{equation}
\end{theorem}

\begin{proof}
  Recall PP $= 4\pi A / P^2$ and $S = P / (2\sqrt{\pi A})$.
  Then:
  \begin{align}
    \frac{1}{\sqrt{\text{PP}}} &= \frac{1}{\sqrt{4\pi A / P^2}} \\
    &= \frac{1}{\sqrt{4\pi A} / P} \\
    &= \frac{P}{\sqrt{4\pi A}} \\
    &= \frac{P}{2\sqrt{\pi A}} = S. \qed
  \end{align}
\end{proof}

\begin{corollary}
  $S(D) \cdot \sqrt{\text{PP}(D)} = 1$ for any polygon $D$.
\end{corollary}

\begin{corollary}
  The monotone transformation $f(x) = 1/\sqrt{x}$ converts PP to S losslessly:
  $S = f(\text{PP})$ and PP $= f(S)^{-1} = 1/S^2$.
\end{corollary}

\subsection{Numerical Verification}

The equivalence $|S - 1/\sqrt{\text{PP}}| < 10^{-9}$ must hold for all
computed districts.
The \textsc{bisect} test suite includes an L0 invariant that verifies this
identity to floating-point precision for all district geometries.

\subsection{Implications}

Theorem~\ref{thm:equivalence} has two practical implications:
\begin{enumerate}
  \item \textbf{No independent computation needed}: If \textsc{bisect} reports PP,
    S is obtained immediately as $S = 1/\sqrt{\text{PP}}$.
    There is no need to run a separate Schwartzberg computation.
  \item \textbf{Statutory compliance}: If a state statute mandates a
    Schwartzberg score threshold (e.g., $S \leq 2.0$), this is equivalent to
    a PP threshold ($\text{PP} \geq 0.25$).
    The \textsc{bisect} label-analyze output can be converted directly.
\end{enumerate}
